\documentclass[10pt,a4paper]{amsart}
\usepackage[margin=19mm]{geometry}
\usepackage{amsmath,amssymb,mathtools}
\usepackage[scr=rsfs,cal=euler]{mathalfa}
\usepackage{xfrac}
\usepackage[unicode,hidelinks,bookmarks=false]{hyperref}
\newcommand{\ie}{i.e.\ }
\newcommand{\cf}{cf.\ }
\newcommand{\resp}{respectively\ }
\newcommand{\Subsection}{\subsection}
\providecommand{\nobreakdash}{\nobreak\hbox{-}}
\setlength{\emergencystretch}{2em}
\multlinegap0pt
\usepackage{bm}
\usepackage{bbm}
\usepackage{dsfont}
\usepackage{xspace}
\usepackage{cancel}
\usepackage{mathtools}
\usepackage{amsthm}
\makeatletter
\@ifundefined{thm}{\newtheorem{thm}{Thm}}{}
\@ifundefined{prop}{\newtheorem{prop}[thm]{Proposition}}{}
\makeatother
\makeatletter
\newcommand{\A}{\mathbb{A}}
\expandafter\ifx\csname AMS\endcsname\relax
\newenvironment{AMS}{}{}
\fi
\DeclareMathOperator{\C}{\mathbb{C}}
\expandafter\ifx\csname keywords\endcsname\relax
\newenvironment{keywords}{}{}
\fi
\newcommand{\s}{{Section}}
\makeatother
\title[On divisibility relation graphs]{On divisibility relation graphs}
\author{Jonathan L. Merzel, J\'an Min\'a\v{c}, Tung T. Nguyen, Nguyen Duy T\^an}
\date{Source: arXiv:2507.06873v1 (CC BY 4.0)}
\begin{document}
\maketitle

\noindent\emph{Source and licence.} On divisibility relation graphs. Version arXiv:2507.06873v1 (CC BY 4.0), distributed under the Creative Commons Attribution 4.0 International licence, as recorded in the arXiv licence metadata for this exact version and shown on the abstract page https://arxiv.org/abs/2507.06873v1. Full rights notice: licenses/arxiv-2507.06873-CC-BY-4.0.txt.\par\smallskip

\noindent\textit{Editorial scope.} Counted unit: the Prop whose source label is \texttt{None} in the article (source0.tex lines 1415--1417, source label \texttt{None}), delivered together with the article's own complete original proof of it (source0.tex lines 1418--1644, one \texttt{proof} environment of 227 lines). No mathematics is written, repaired or re-derived: statement and proof are the article's own text line for line. Required context retained uncounted: none. Every definition, display and label of those lines is retained unchanged. Omitted from this package and not redistributed: 1--1414, 1645--1782. Nothing inside a retained excerpt is omitted or altered: every retained body is delivered exactly as the article's lines are written, so no omission receipt of any kind accompanies this package. Citations: the retained excerpts carry no citation command at all, so no bibliography entry of the article is transcribed and no reference is redistributed. Reference adaptation: none was needed and none was made; every reference inside the delivered unit resolves inside it, so no unresolved reference or citation is produced. Printed slips kept literal: no printed word, symbol, hypothesis or number of the counted statement or of its proof was altered. Layout and class: the article's own class is \texttt{amsart} with packages including amsmath, amssymb, amsfonts, float, xy, caption, enumerate, mathpazo, a4wide, tikz, soul, bm, graphicx, hyperref; the wrapper is the lane's pinned \texttt{amsart} editorial wrapper at 10pt on A4, which restates the theorem-like environments the retained text needs and adds the article's own macro definitions that the retained text needs (\texttt{\textbackslash A}, \texttt{\textbackslash C}, \texttt{\textbackslash s}), with extra wrapper packages (bm, bbm, dsfont, xspace, cancel, mathtools). The delivered numbering is the unit's own. Assets: no retained span contains a figure environment, a graphics-inclusion command, a PDF-page inclusion, a table or a TikZ picture; no image file is copied into this package, the package declares no asset dependency and no image file is redistributed with it. Licence: version arXiv:2507.06873v1 of this article is distributed under the Creative Commons Attribution 4.0 International licence, as recorded in the arXiv licence metadata for this exact version and shown on the abstract page https://arxiv.org/abs/2507.06873v1. A full rights notice is delivered at licenses/arxiv-2507.06873-CC-BY-4.0.txt. Characters: every retained excerpt is pure ASCII, so the delivered engine, XeLaTeX with its default Latin Modern text font, reports no missing character.
\begin{prop}
If $n=p^a q^b$ and $a \equiv b \equiv 1 \pmod{6}$ then $0$ is an eigenvalue of $D_n.$
\end{prop}

\begin{proof} Suppose that     \[
n = p^u q^v \quad \text{with} \quad u \equiv v \equiv 1 \pmod{6}.
\]
List the divisors of \( n \) in the order:
\[
1, q, q^2, \ldots, q^v, p, pq, pq^2, \ldots, pq^v, \ldots, p^u, p^u q, \ldots, p^u q^v.
\]
The adjacency matrix of \( D_n \) is of the form:
\[
M = 
\begin{bmatrix}
V & U & U & \cdots & U \\
U^T & V & U & \cdots & U \\
U^T & U^T & V & \cdots & U \\
\vdots & \vdots & \vdots & \ddots & \vdots \\
U^T & U^T & U^T & \cdots & V
\end{bmatrix},
\]
where
\[
V = 
\begin{bmatrix}
0 & 1 & 1 & \cdots & 1 \\
1 & 0 & 1 & \cdots & 1 \\
1 & 1 & 0 & \cdots & 1 \\
\vdots & \vdots & \vdots & \ddots & \vdots \\
1 & 1 & 1 & \cdots & 0 
\end{bmatrix}_{(v+1)\times(v+1)}
\quad \text{and} \quad
U = 
\begin{bmatrix}
1 & 1 & \cdots & 1 \\
0 & 1 & \cdots & 1 \\
\vdots & \vdots & \ddots & \vdots \\
0 & 0 & \cdots & 1
\end{bmatrix}_{(v+1)\times(v+1)}
\]
\setcounter{MaxMatrixCols}{20}
Consider the following column matrices of size \((v+1) \times 1\):
\[
\begin{aligned}
    A&=\left[\begin{array}{*{15}c}
            0& 1&1& 0 &-1&-1  &\cdots &0 & 1&1& 0&-1&-1 & 0 & 1 \\
           \end{array}\right]^T\\
    A'&=\left[\begin{array}{*{15}c}
           -1& 0& 1&1& 0 &-1  &\cdots &-1&0 & 1&1& 0&-1&-1 & 0  \\
           \end{array}\right]^T.
\end{aligned}
\]
(In $A$, the block $\begin{bmatrix}
     0& 1&1& 0 &-1&-1 
\end{bmatrix}$ is repeated $(v-1)/6$ times. The matrix $A'$ is obtained from $A$ by left-shifting.) Similary, we consider following column matrices of size \((v+1) \times 1\):
\[
\begin{aligned}
       B&=\left[\begin{array}{*{15}c}
           -1& -1& 0& 1 &1 &0 &\cdots &-1& -1& 0& 1 &1 &0 & -1 & -1  \\
           \end{array}\right]^T,\\
        B'&=\left[\begin{array}{*{15}c}
            0& -1&-1& 0 &1&1  &\cdots &0 & -1&-1& 0&1&1 & 0 & -1 \\
           \end{array}\right]^T\,\\
\end{aligned}
\]
and
\[
\begin{aligned}
       C&=\left[\begin{array}{*{15}c}
           1& 0& -1& -1 &0 &1 &\cdots &1& 0& -1& -1 &0 &1 &1 &0  \\
           \end{array}\right]^T,\\
        C'&=\left[\begin{array}{*{15}c}
            1& 1&0& -1 &-1&0  &\cdots &1& 1&0& -1 &-1&0&1&1  \\
           \end{array}\right]^T.
\end{aligned}
\]
Let $X$ be the following column matrix of size $[(u+1)(v+1)]\times 1$:
\[
X=\begin{bmatrix}
A & A' & B & B' & C&C' &\cdots & A & A' & B & B' & C&C' & A & A'
\end{bmatrix}^T.
\]
(In $X$, the block $\begin{bmatrix}
   A & A' & B & B' & C&C' 
\end{bmatrix}$ is repeated $(u-1)/6$ times.) 


We claim  $MX=0$. In particular, it shows that $0$ is an eigenvalue of $M$.

We need to show that for every non-negative integers $s,t$ with $s+t=u\equiv 1\pmod 6$, one has
\[
\begin{bmatrix}
    U^T & \cdots &U^T & V &U&\cdots & U
\end{bmatrix} X =0.
\]
Here in the first matrix,  $U^T$ appears $s$ times and $U$ appears $t$ times.

Note that $A+A'+B+B'+C+C'=0$. Hence $U^T(A+A'+B+B'+C+C')=0=U(B+B'+C+C'+A+A')$. We only need to consider the following six cases.

\medskip
\noindent {\bf Case 1}: $s\equiv 0\pmod 6$ and $t\equiv 1\pmod 6$. In this case, it suffices to show that \[\begin{bmatrix}
    V & U 
\end{bmatrix} \begin{bmatrix}
    A\\A'
\end{bmatrix}=0.\]

Let $\begin{bmatrix}
    V & U 
\end{bmatrix} \begin{bmatrix}
    A\\A'
\end{bmatrix}=\begin{bmatrix}s_1\\s_2\\\vdots\end{bmatrix}$. We have
\[
s_1=\sum_{i\geq 2}a_i+ \sum_{i\geq 1}a'_i= 1+(-1)=0.
\]
On the other hand, the difference between the $(i+1)$-st row and the $i$th row of $[V \mid U]$ is the row
\[
\left [\; 0\quad \cdots\quad 0 \quad 1 \quad -1 \quad 0\quad\cdots \quad 0\quad \mid\quad  0\quad\cdots \quad 0\quad -1\quad 0\quad\cdots \quad 0\;\right].
\]
Here, the first $-1$ is in the $i+1$-th position and the second $-1$ is in the $(s+1)+i$-th position. We deduce that
\[
s_{i+1}-s_i=a_i-a_{i+1}-a'_i=0.
\]
Hence $s_i=0$, for every $i$.


\medskip
\noindent {\bf Case 2}: $s\equiv 1\pmod 6$ and $t\equiv 0\pmod 6$. In this case, it suffices to show that \[\begin{bmatrix}
    U^T & V 
\end{bmatrix} \begin{bmatrix}
    A\\A'
\end{bmatrix}=0.\]

Let $\begin{bmatrix}
    U^T & V 
\end{bmatrix} \begin{bmatrix}
    A\\A'
\end{bmatrix}=\begin{bmatrix}s_1\\s_2\\\vdots\end{bmatrix}$. We have
\[
s_1=a_1+ \sum_{i\geq 2}a'_i= 0+0=0.
\]
On the other hand, the difference between the $(i+1)$-st row and the $i$th row of $[U^T \mid V]$ is the row
\[
\left [\; 0\quad \cdots\quad 0 \quad 1  \quad 0\quad\cdots \quad 0\quad \mid\quad  0\quad\cdots \quad 0 \quad  1\quad -1\quad 0\quad\cdots \quad 0\;\right].
\]
Here, the first $1$ is in the $i+1$-th position and the second $1$ is in the $(s+1)+i$-th position. We deduce that
\[
s_{i+1}-s_i=a_{i+1}+a'_i-a'_{i+1}=0.
\]
Hence $s_i=0$, for every $i$.

Using  similar arguments as in the previous two cases, we can show that
\[
VB+UB'= [\;s \quad s\quad\cdots\quad s\;]^T= U^TB+VB',
\]
where $s=\sum_{i\geq 2}b_i+ \sum_{i\geq 1}b'_i=-2=b_1+ \sum_{i\geq 2}b'_i$, and that
\[
\begin{aligned}
&VA'+UB=U^TA'+VB=[-2\quad -2\quad \cdots \quad -2\;]^T,\\
&VB'+UC= U^TB'+VC=0, \\
&VC+UC'=U^TC+VC'=[\;2\quad 2\quad \cdots \quad 2\;]^T.   
\end{aligned}
\]



\medskip
\noindent {\bf Case 3}: $s\equiv 2\pmod 6$ and $t\equiv 5\pmod 6$. In this case, it suffices to show that \[\begin{bmatrix}
    U^T & U^T & V & U & U & U&U&U 
\end{bmatrix} \begin{bmatrix}
    A\\A'\\B\\B'\\C\\C'\\A\\A'
\end{bmatrix}=0.\]

We have
\[
\begin{aligned}
   &U^T(A+A')+VB+U(B'+C+C'+A+A')=-VA'+U^TA' +VB-UB \\
   &= (U^TA'+VB)-(VA'+UB)=0.
\end{aligned}
\]

\medskip
\noindent {\bf Case 4}: $s\equiv 3\pmod 6$ and $t\equiv 4\pmod 6$. In this case, it suffices to show that \[\begin{bmatrix}
    U^T & U^T & U^T & V & U & U & U&U 
\end{bmatrix} \begin{bmatrix}
    A\\A'\\B\\B'\\C\\C'\\A\\A'
\end{bmatrix}=0.\]

We have 
\[
\begin{aligned}
&U^T(A+A'+B) + VB'+U(C+C'+A+A')\\
&=U^T(A+A')+VB+U(B'+C+C'+A+A') +U^TB+VB' -VB -UB' \\
&=U^TB+VB' -(VB +UB')=0.
\end{aligned}
\]


\medskip
\noindent {\bf Case 5}: $s\equiv 4\pmod 6$ and $t\equiv 3\pmod 6$. In this case, it suffices to show that \[\begin{bmatrix}
    U^T & U^T & U^T & U^T&V & U & U & U 
\end{bmatrix} \begin{bmatrix}
    A\\A'\\B\\B'\\C\\C'\\A\\A'
\end{bmatrix}=0.\]

We have 
\[
\begin{aligned}
&U^T(A+A'+B+B') + VC+U(C'+A+A')\\
&=U^T(A+A'+B) + VB'+U(C+C'+A+A') +U^TB'+VC -VB' -UC \\
&=U^TB'+VC  -(VB' +UC)=0.
\end{aligned}
\]


\medskip
\noindent {\bf Case 6}: $s\equiv 5\pmod 6$ and $t\equiv 2\pmod 6$. In this case, it suffices to show that \[\begin{bmatrix}
    U^T & U^T & U^T &U^T &U^T & V & U & U 
\end{bmatrix} \begin{bmatrix}
    A\\A'\\B\\B'\\C\\C'\\A\\A'
\end{bmatrix}=0.\]
We have 
\[
\begin{aligned}
&U^T(A+A'+B+B'+C) + VC'+U(A+A')\\
&=U^T(A+A'+B+B') + VC+U(C'+A+A')+U^TC+VC' - VC -UC' \\
&=U^TC+VC'  -(VC +UC')=0.
\end{aligned}
\]

\end{proof}

\end{document}
